\documentclass[aspectratio=169]{beamer}
\usetheme{metropolis}
\usepackage{graphicx}
\usepackage{svg}
\usepackage{amsmath}
\title{MOST, Stable Boundary Layer, and Curvature-aware Corrections}
\author{Your Name}
\date{\today}

\begin{document}
\maketitle

\begin{frame}{Motivation}
  \begin{itemize}
    \item \textbf{Problem:} Coarse vertical grids misclassify stability, yielding warm biases and pollutant errors.
    \item \textbf{Goal:} Resolution-aware correction that preserves neutral physics and reduces bias.
    \item \textbf{Scope:} Curvature analysis for the \textbf{stable boundary layer (SBL)}, not BD USL parameters.
  \end{itemize}
  \note{
    Set the stage: stability misclassification near the surface has outsized impacts.
    Emphasize that our curvature discussion is SBL-specific. BD USL applies to unstable daytime turbulence and is used only for contrast.
  }
\end{frame}

\begin{frame}{MOST essentials}
  \begin{itemize}
    \item \textbf{Key variable:} \(\zeta = z/L\) (dimensionless height).
    \item \textbf{Functions:} \(\phi_m(\zeta), \phi_h(\zeta)\) govern shear/heat transfer.
    \item \textbf{Related measures:} \(Ri_g\), stability via shear and buoyancy.
  \end{itemize}
  \note{
    Keep intuition: near-surface turbulence collapses to \(\zeta\).
    We'll link \(\phi_m, \phi_h\) to curvature via a compact fingerprint \(\Delta\).
  }
\end{frame}

\begin{frame}{Curvature fingerprints: SBL vs BD USL}
  \begin{center}
    \includesvg[width=0.9\linewidth]{curvature_sbl_usl.svg}
  \end{center}
  \begin{itemize}
    \item \textbf{SBL:} Concave-down Ri\(_g\) vs \(\zeta\) (negative curvature; \(\Delta < 0\)).
    \item \textbf{BD USL:} Unstable branch often nearer neutral/concave-up; different parameter regime.
    \item \textbf{Clarification:} Our correction targets SBL curvature only.
  \end{itemize}
  \note{
    Point at the blue curve: SBL concave-down behavior.
    Red dashed is BD USL—showing contrast but not used for parameterization here.
  }
\end{frame}

\begin{frame}{Neutral curvature invariant}
  \begin{itemize}
    \item \textbf{Fingerprint:} \(\Delta = \alpha_h \beta_h - 2\,\alpha_m \beta_m\).
    \item \textbf{Interpretation:} Sign of \(\Delta\) sets concavity of stability curves.
    \item \textbf{SBL fits:} Typically yield \(\Delta < 0\) (concave-down).
  \end{itemize}
  \note{
    You can motivate \(\Delta\) as a compact way to preserve neutral behavior while correcting curvature effects.
    No need to deep-derive; focus on physical meaning.
  }
\end{frame}

\begin{frame}{Layer averaging bias in SBL}
  \begin{center}
    \includesvg[width=0.9\linewidth]{grid_bias.svg}
  \end{center}
  \begin{itemize}
    \item \textbf{Issue:} Averaging across a concave-down profile yields \(Ri_b < Ri_g(z_g)\).
    \item \textbf{Geometric mean height:} \(z_g = \sqrt{z_0 z_1}\) is the natural evaluation point.
    \item \textbf{Impact:} Systematic underestimation of stability on coarse grids.
  \end{itemize}
  \note{
    Invoke Jensen’s inequality informally: concave-down -> average below the curve.
    Emphasize evaluating at \(z_g\) rather than linear midpoints.
  }
\end{frame}

\begin{frame}{Correction strategy}
  \begin{itemize}
    \item \textbf{Preserve:} Neutral behavior and the curvature fingerprint \(2\Delta\).
    \item \textbf{Adjust:} Bulk Ri\(_b\) toward the true Ri\(_g(z_g)\) with grid-aware damping.
    \item \textbf{Converge:} Correction vanishes as \(\Delta z \to 0\).
  \end{itemize}
  \note{
    Core principle: don’t break neutrality; correct only the curvature-induced bias.
    Stress resolution-awareness and graceful convergence.
  }
\end{frame}

\begin{frame}{Validation scaffold}
  \begin{center}
    \includesvg[width=0.9\linewidth]{validation_layout.svg}
  \end{center}
  \begin{itemize}
    \item \textbf{Comparison:} Observations vs model (uncorrected) vs corrected.
    \item \textbf{Metrics:} RMSE reduction indicative of improved stability depiction.
    \item \textbf{Cases:} Apply across diverse SBL nights to test robustness.
  \end{itemize}
  \note{
    Replace placeholders with your tower/time-series data.
    Narrate RMSE changes and qualitative fit improvements (phase, amplitude).
  }
\end{frame}

\begin{frame}{Limitations and open questions}
  \begin{itemize}
    \item \textbf{Transferability:} Parameter portability across terrains and urban canopies.
    \item \textbf{Inflections:} Handling mixed regimes and nocturnal jets.
    \item \textbf{L(z):} Height-varying Obukhov length effects.
  \end{itemize}
  \note{
    Keep it candid: where the approach might struggle.
    Invite collaboration on canopy and complex terrain extensions.
  }
\end{frame}

\begin{frame}{Takeaways}
  \begin{itemize}
    \item \textbf{Curvature matters:} It sets departures from neutrality.
    \item \textbf{Evaluate smartly:} Use \(z_g\) and preserve \(2\Delta\).
    \item \textbf{Correct gently:} Grid-aware adjustments reduce biases without breaking physics.
  \end{itemize}
  \note{
    Summarize and reiterate SBL-only applicability. BD USL parameters are not the source of the curvature used in this correction.
  }
\end{frame}

\end{document}
